The analysis contains 2110 measurements across five conditions, with group sizes ranging from 319 to 566. Attenuation spans 10.0--30.0 dB, QBER spans 0.0261--0.0865, and raw SKR spans 0--354598 bits/s. There are 277 zero-SKR measurements (13.1\%); these are retained, and the $\log_{10}(\mathrm{SKR}+1)$ transformation preserves their value at zero.

The strongest absolute pairwise correlation is between LogSKR and QBER ($r=-0.957$). In particular, attenuation correlates 0.878 with QBER and -0.824 with transformed SKR. These pooled associations combine within-condition trends and differences in the sampled condition/attenuation mix; they do not by themselves isolate a causal effect of injected power.

For example, mean sampled attenuation is 23.13 dB for baseline and 19.61 dB for 12 dBm, while their median raw SKR values are 3541 and 91610 bits/s, respectively. Those medians summarize different attenuation distributions. A raw median comparison must therefore not be interpreted as a controlled injected-power effect.


At the shared attenuation of 28 dB, mean SKR is 2969 bits/s for baseline and 0 bits/s for 12 dBm. The 12 dBm condition also has 22.1\% zero-throughput readings overall. This matched-setting comparison explains why its high pooled median is not evidence of better performance under stronger injection.
